% ==============================================================================
% Section: Conclusion
% Status: STYLE EDIT (2026-04-07 — removed bold labels, flowing prose)
% ==============================================================================

\section{Conclusion}
\label{sec:conclusion}

Three findings emerge from the calibrated SFE analysis. First, the settlement
cap is binding under concentrated supply: with $K \leq 3$ pivotal suppliers,
the SFE equilibrium for 2026/27 falls at the VRR Point~(a) cap of
\$329.17/MW-day, exceeding the settlement cap $\bar{p}_S = \$325$ by
\$4.17/MW-day at the RTO level and by \$115--\$141/MW-day in constrained
LDAs. Second, the \$21~billion savings figure depends on the counterfactual:
using the SFE equilibrium rather than the VRR Point~(a) price yields a
two-year RTO transfer of approximately \$612~million, substantially smaller
than the stated figure; both magnitudes are sensitive to which Point~(a) the
pre-settlement VRR is taken to fix, and the at-cap outcome in 2026/27 and
2027/28---each of which cleared below its reliability requirement---is
consistent with both the SFE mechanism and physical scarcity
(Sections~\ref{sec:cap_incidence} and~\ref{subsec:case_2027}). Third, structural channels can reduce prices without truncating the
signal: CETL expansion through PJM's Regional Transmission Expansion
Plan \citep{PJM_RTEP}, VRR slope recalibration that disciplines supply
at every price below the cap rather than via direct truncation
\citep{FERC2025_ER25-1357}, and supply-side intervention via
interconnection-queue reform and expanded demand response and storage
participation \citep{FERC2023_order2023, FERC2024_order2023A} each
lower the SFE equilibrium through changes in market structure or
demand-curve shape rather than through price truncation.

The settlement cap's allocative significance depends on its proximity
to the investment threshold. In 2026/27, the cap of \$325/MW-day exceeds
RTO Net CONE (\$212.14) by 53\%, so the cap alone does not prevent
new entry from being financially viable. The more substantive concern
is informational: because the cap binds in consecutive years, potential
entrants observe the cap rather than the unconstrained equilibrium price
and cannot distinguish a market that would have cleared at \$333 from
one that would have cleared far above \$325, reducing the
informativeness of the capacity price for investment decisions
\citep{Hogan2005_energy_adequacy, Joskow2008_capacity_payments}---a
concern amplified in constrained LDAs where local Net CONE values
exceed the settlement cap (e.g., PS and PS NORTH Net CONE
$= \$369$/MW-day in 2026/27).

The model's contribution is evaluative rather than prescriptive. A temporary
cap can reduce payments, but the comparative statics show that the
underlying concentration and transmission conditions generating high
equilibrium prices are what place the equilibrium at the cap in the first
place. The structural channels identified in the comparative statics
(CETL expansion, VRR slope, entry) would reduce equilibrium markups through
market structure rather than by truncating the price signal; the model does
not rank reforms or quantify their welfare consequences, and whether any
particular change is warranted depends on regulatory objectives and
cost--benefit considerations outside its scope. Natural extensions include an
asymmetric-firm SFE, a dynamic model connecting markups to capacity adequacy,
and an empirical structural model recovering firm-level costs from offer data.

